% What's expressible as native Quarto/pandoc PDF metadata lives in
% _extension.yml instead (documentclass, classoption, geometry,
% colorlinks/linkcolor/urlcolor/citecolor). Only what has no YAML
% equivalent stays as raw LaTeX here.

% fonts.latex is overridden as a no-op (see partials/pdf/fonts.latex), so
% fontenc/inputenc need loading explicitly -- \fontencoding{T1} (small-caps
% fix below) needs a T1 encoding to select.
\usepackage[utf8]{inputenc}
\usepackage[T1]{fontenc}

% Language: driven by Quarto's own native `lang:` document metadata (not
% a custom key of this extension's) -- Quarto's default template already
% turns `lang: en`/`lang: fr` into a `babel-lang` documentclass option
% and loads babel accordingly (see doc-class.tex/babel-lang.tex in
% Quarto's own PDF template), so babel is already active with the right
% language by the time this file is spliced in. \ifenglish just exposes
% that same choice to consuming documents' own raw LaTeX (e.g.
% personal-macros.tex's EN/FR vocabulary macros) via babel's own
% \iflanguage test -- no manual classoption/`\selectlanguage` needed.
\newif\ifenglish
\iflanguage{english}{\englishtrue}{\englishfalse}

% Header/footer: no native Quarto equivalent.
\usepackage{fancyhdr}
\pagestyle{fancy}
\fancyhf{}
\fancyhead[R]{}
\fancyfoot[RO,LE]{\thepage}
\fancyfoot[LO,RE]{\footnotesize\sc\footer}
\renewcommand{\headrulewidth}{0pt}
\renewcommand{\footrulewidth}{0pt}

% Section titles: package load only -- \titleformat calls are in
% before-body.tex, alongside the colors they use.
\usepackage[explicit]{titlesec}

% array: needed for the custom R{<width>} column type \multblock relies
% on. Also loaded by Quarto's own tables.tex, but load it explicitly so
% \multblock keeps working if a document ever disables that (`tables: false`).
\usepackage{array}

\usepackage{url}
\usepackage{eurosym}
\usepackage{fontawesome}

% float: needed for the [H] placement specifier, so a captioned figure
% (e.g. in a "Research Project" section from executable code) stays
% exactly where it's placed instead of drifting as a normal LaTeX float
% -- see the extension README, "Executable code". Not needed for tables:
% pandoc emits captioned tables as `longtable`, which isn't a float.
\usepackage{float}

% This document's own personal macros (\nom, \title, ...) and any EN/FR
% vocabulary it needs -- see the extension README, "Usage".
